\section{Recovery for every fixed \texorpdfstring{$K$}{K} at general density}
\label{sec:main-results}

We first state the complete algorithmic guarantee, then identify the deterministic replacement argument and the probability estimates that prove it. Throughout, the model satisfies \eqref{eq:model-conditions}--\eqref{eq:density-condition} and the chosen $c_0$ satisfies the admissibility condition following \eqref{eq:profiles}.

\subsection{The oracle information and the main theorem}
For communities $a\ne b$, define the exact product-Bernoulli Chernoff information, using Zhou--Li's profile notation,
\begin{align}
 D_\alpha(p_{a,-i}\Vert p_{b,-i})
  &=-\sum_{j\ne i}\log\{p_{aj}^{1-\alpha}p_{bj}^{\alpha}
           +(1-p_{aj})^{1-\alpha}(1-p_{bj})^{\alpha}\},
  \label{eq:chernoff-information}\\
 D^*_{i,ab}&=\max_{\alpha\in[0,1]}
                 D_\alpha(p_{a,-i}\Vert p_{b,-i}),
 &I_n&=\min_{i,\ a\ne b}D^*_{i,ab}.
 \label{eq:exact-information}
\end{align}
The absent self-loop coordinate is excluded. The uniform shape and separation conditions imply $I_n\asymp np$, even if $P/p$ does not converge. For the binary experiment that reveals the other labels and the true parameters, with equal prior probabilities or fixed priors bounded away from zero and one, the oracle Bayes risk satisfies
\begin{equation}
 R_{i,ab}^{\mathrm{or}}=\exp\{-(1+o(1))D^*_{i,ab}\}.
 \label{eq:oracle-risk}
\end{equation}
The appendix proves this leading form directly by exponential tilting. Zhou and Li~\citep{zhouli2020} studied optimal error rates; no refined multiplicative rate from that work is asserted here.

\begin{theorem}[Nearly optimal spectral likelihood recovery]
\label{thm:main-recovery}
For every fixed $K\ge2$, Algorithm~1 obeys
\begin{equation}
 \E\Mis(\widehat z,z)\le\exp\{-(1+o(1))I_n\}.
 \label{eq:main-rate}
\end{equation}
If $\liminf_{n\to\infty}I_n/\log n>1$, then
$\Pr\{\Mis(\widehat z,z)>0\}\to0$.
The initialization and replacement use only the retained eigenpairs, and no weak-initialization assumption is imposed. Growing and replacement trials may use any finite number of exact EM updates, with at least one M-step per fit. The mandatory final E/M update is part of the algorithm. Any further finite sequence of exact EM updates, including a data-dependent stopping time, preserves the bound. With polynomial iteration caps, the algorithm takes polynomial time.
\end{theorem}

The density condition is $p=\Omega(\log n/n)$. Together with probability comparability, it implies expected degrees $\Omega(\log n)$ and permits larger orders. This theorem covers both sparse sequences $p\to0$ and dense sequences with probabilities bounded away from one. The proof does not require a fixed normalized matrix or a restriction such as $np=O(\log n)$.

The guarantee attains the leading binary oracle Chernoff exponent. Revealing the other labels and parameters can only help binary testing, which explains the benchmark. This argument is not a new global minimax lower bound for an unspecified parameter class.

\subsection{The deterministic replacement contraction}
Define the profile divergence and the hard and soft fitting losses by
\begin{align}
 D(x\Vert y)&=\sum_j\left\{x_j\log\frac{x_j}{y_j}
                  +(1-x_j)\log\frac{1-x_j}{1-y_j}\right\},
 \label{eq:profile-divergence}\\
 G(\mathcal C)&=\sum_i\min_kD(q_i\Vert\mu_k),
 &\mathcal C&=(\mu_1,\ldots,\mu_K),\notag\\
 F(\widehat\pi,\mathcal C)
    &=-\sum_i\log\sum_k\widehat\pi_k e^{-D(q_i\Vert\mu_k)}.
 \label{eq:soft-loss}
\end{align}
There is the exact identity
\begin{equation}
 \mathcal L=S_q-F,\quad S_q=\sum_i\ell_i(q_i),\qquad
 G\le F,\qquad F(K^{-1}\mathbf1,\mathcal C)\le G+n\log K.
 \label{eq:soft-hard}
\end{equation}
Thus selecting the greatest observed working likelihood minimizes $F$.

The population argument uses exactly $K$ true profiles and $K$ fitted profiles. Put
$G_*(\mathcal C)=\sum_an_a\min_kD(p_{a*}\Vert\mu_k)$.

\begin{lemma}[A replacement removes at least a $1/K$ fraction]
\label{lem:population-replacement}
For any list $\mathcal C$ of $K$ fitted profiles, there exist $a,r\in[K]$ such that replacing $\mu_r$ by $p_{a*}$ gives
\begin{equation}
 G_*(\mathcal C^{r\leftarrow p_{a*}})
       \le(1-1/K)G_*(\mathcal C).
 \label{eq:population-contraction}
\end{equation}
Duplicate fitted profiles are allowed. The conclusion uses only nonnegative losses and $D(p_{a*}\Vert p_{a*})=0$.
\end{lemma}

\begin{proof}
For every true class choose a minimizing center $\phi(a)$. Choose $a$ whose contribution $n_aD(p_{a*}\Vert\mu_{\phi(a)})$ is greatest, and hence at least $G_*/K$. If $\phi:[K]\to[K]$ is not onto, delete an index outside its image. If it is onto, it is a bijection; delete $\phi(a)$. In both cases all other classes retain a chosen old center. Inserting $p_{a*}$ sets the chosen class's loss to zero without increasing any other class's loss.
\end{proof}

The algorithm does not know $\phi$ or $p_{a*}$. It tests all deletion positions and uses observed node candidates. The uniform average approximation below makes this population argument stable for those operations.

\begin{theorem}[Deterministic likelihood-loss contraction]
\label{thm:replacement-contraction}
Suppose the observed and true profiles, and the initial fitted means, lie in
\[
 [c p,\,\min\{C p,1-b\}]^n
\]
for fixed positive $c,C,b$, and
\begin{equation}
 \frac1{n^2p}\sum_i\|q_i-p_{z_i,*}\|_1\le a_n\to0.
 \label{eq:replacement-input}
\end{equation}
Assume $n_a\ge\pi_0n$, $np\to\infty$, and fixed $K$. The replacement rounds in Algorithm~1 satisfy, for constants $B,M$ independent of the round,
\begin{align}
 F_{t+1}&\le (1-1/K)F_t+B a_n n^2p+n\log K,
 \label{eq:replacement-contraction}\\
 \frac{F_{T_n}}{n^2p}
 &\le M(1-1/K)^{T_n}+KB a_n+\frac{K\log K}{np}
   =o(1).
 \label{eq:replacement-weak-loss}
\end{align}
This holds for any initial simplex weights, including zero weights, and does not assume convergence of any trial to a stationary point or global optimum.
\end{theorem}

On the common box, $D$ is uniformly Lipschitz in each profile's $\ell_1$ norm. Therefore $|G-G_*|\le C a_n n^2p$ simultaneously for every possible fitted center list. Each true class has an observed candidate within $(a_n/\pi_0)np$ of its true profile. The best empirical replacement consequently inherits \eqref{eq:population-contraction} up to $O(a_n n^2p)$. Equation~\eqref{eq:soft-hard}, uniform trial weights and EM monotonicity convert this to \eqref{eq:replacement-contraction}. The proof in Appendix~\ref{app:replacement-proof} supplies all perturbation details.

A single replacement pass only gives a fixed contraction and does not establish \eqref{eq:replacement-weak-loss}. Repetition is the algorithmic step that removes the general-$K$ initialization gap. In contrast, the complete growing-only trajectory need not preserve a pure component at every stage; the proof does not assume that property.

\subsection{Full-graph spectral approximation with sufficient confidence}
Let $\Pi_*=Z(Z^\top Z)^{-1}Z^\top$ and define the comparison
\[
 q^0=\operatorname{clip}(A\Pi_*,\varepsilon_n,1-\varepsilon_n).
\]
The true labels and separate leave-one-out eigendecompositions are not algorithm inputs.

\begin{lemma}[Full-graph approximation at general density]
\label{lem:full-spectral}
For every fixed $H>0$, there are $C_H<\infty$, a deterministic $a_{H,n}\to0$, and an event $G_H$ with
\begin{equation}
 \Pr(G_H^c)\le C_H e^{-Hnp},
 \label{eq:spectral-failure}
\end{equation}
on which
\begin{align}
 \max_i\|q_i-q_i^0\|_1&\le a_{H,n}np,
 &\frac1{n^2p}\sum_i\|q_i-p_{z_i,*}\|_1&\le a_{H,n},
 \label{eq:general-profile-errors}\\
 cp\le q_{ij}&\le\min\{C_Hp,1-cp\},
 &\frac\kappa4np\le s&\le2np.
 \label{eq:general-profile-box}
\end{align}
The constants depend on the fixed model bounds and the chosen $c_0$. In particular all true probabilities lie inside the clipping interval on $G_H$ for sufficiently large $n$.
\end{lemma}

The probability in \eqref{eq:spectral-failure} is needed when $I_n\gg\log n$. An arbitrary fixed inverse-polynomial failure bound need not be smaller than $e^{-I_n}$. We instead use the exponential tail of Bandeira and van Handel~\citep[Corollary 3.12 and Remark 3.13]{bandeiravhandel2016}, and verify the row concentration in Abbe et al.~\citep[Theorem 2.1 and Lemma 7]{abbe2020} at that same scale. Appendix~\ref{app:general-full-spectral} treats the signed signal subspaces and the exact no-self-loop mean.

The profile box in \eqref{eq:general-profile-box} is uniformly bounded away from one as well: when $C_Hp\le1/2$ this is immediate, and otherwise $cp\ge c/(2C_H)$. It therefore supplies the hypothesis of Theorem~\ref{thm:replacement-contraction}.

\subsection{Low fitting loss and the final likelihood update}
\begin{lemma}[Small soft loss certifies weak responsibilities]
\label{lem:low-loss-responsibilities}
For every deterministic $\xi_n\to0$, on $G_H$ consider all fits with exactly $K$ profiles in the common box and
$F\le\xi_n n^2p$. After one common label permutation for each fit, its exact E-step responsibilities obey
\begin{equation}
 \frac1n\sum_{i,a}|R_{ia}-Z_{ia}|=o(1).
 \label{eq:weak-responsibilities}
\end{equation}
The error bound is deterministic and uniform over this class of fits and all simplex weights; no independent initialization is required.
\end{lemma}

The proof first uses the strong convexity bound
$D(x\Vert y)\ge\|x-y\|_1^2/(2C_Hnp)$.
Together with the average profile approximation and separation of the true profiles, this matches the $K$ fitted profiles to the $K$ true profiles. The exact variational identity
\begin{equation}
 F_i=\sum_aR_{ia}D(q_i\Vert\widehat p_a)
             +\operatorname{KL}(R_i\Vert\widehat\pi)
       \ge\sum_aR_{ia}D(q_i\Vert\widehat p_a)
 \label{eq:variational-loss}
\end{equation}
then bounds the responsibility mass on wrongly matched profiles. This is why an objective-based selection suffices even if a trial stopped before convergence.

\begin{theorem}[Local empirical EM preserves the oracle exponent]
\label{thm:local-spectral-em}
Let an initial responsibility matrix, possibly using the whole graph, have aligned average error at most a deterministic $\gamma_n\to0$ except on an event of probability $\delta_n$, where
\[
 \liminf_{n\to\infty}\frac{-\log\delta_n}{I_n}\ge1.
\]
One profile and weight M-step on $q$, followed by likelihood classification, satisfies \eqref{eq:main-rate}. The conclusion holds after any further finite exact EM updates and any finite data-dependent stopping time after that first M-step.
\end{theorem}

The M-step gives parameter-profile error $O\{(\gamma_n+a_{H,n})np\}$. For rare rows, clipping is handled directly by a Chernoff moment argument with a deterministic envelope for the random spectral threshold. It costs at most $2^K$ in the moment, including both clipping sides; it is not asserted to perturb every rare-row score by $o(I_n)$. A deterministic invariant-basin argument controls all later updates on the same graph event. These steps are proved in Appendix~\ref{app:general-local-em}.

Combining Theorem~\ref{thm:replacement-contraction}, Lemma~\ref{lem:low-loss-responsibilities}, and Theorem~\ref{thm:local-spectral-em} proves Theorem~\ref{thm:main-recovery}. The only exceptional event is $G_H^c$; choose a fixed $H>1+\limsup I_n/(np)$. No additional randomized seeding failure remains.

\subsection{Why the complete nonedge term matters}
The sparse score used in earlier experiments is
$\sum_jq_{ij}\log\widehat p_j-\sum_j\widehat p_j$.
For $p=o(1)$, its oracle contrast differs from the complete Bernoulli contrast by
$O\{p\deg(i)+np^2\}=o(np)$
on the appropriate degree event. The needed remainder is relative to $I_n$, not an absolute $o(1)$. This permits expected degrees far above $\log n$ within sparse models.

For dense graphs, the discarded term can change the leading exponent. In the balanced model
\[
 P=\begin{pmatrix}0.8&0.2\\0.2&0.2\end{pmatrix},
\]
only edges to the first block discriminate the hypotheses. For $m$ such potential neighbors, the correct threshold is $N/m=1/2$, with information $0.22314355\,m$. The sparse score instead uses $N/m=0.6/\log4$, whose smaller directional error exponent is $0.13905314\,m$. This loss occurs even with the true parameters. Algorithm~1 retains the full Bernoulli expression to cover the full density range of \eqref{eq:density-condition}.
